\phantomsection\addcontentsline{toc}{section}{Smart Cities}
\ECchapter{13}{Smart Cities}{Can connected systems make cities more sustainable?}{ECGreen}

\ECObjectives{
\begin{itemize}
\item Describe the architecture and control loop of a connected urban service.
\item Quantify how sensing, communication and control can reduce resource consumption.
\item Justify choices involving interoperability, privacy, resilience and social inclusion.
\end{itemize}}

\begin{multicols}{2}
\begin{engineeringnews}
Cities are installing connected lighting, traffic and environmental-monitoring systems. The most
successful projects do not simply collect data: they convert selected measurements into useful local
decisions while protecting citizens and keeping essential services available during failures.
\end{engineeringnews}

\begin{ecarticle}
A smart city uses sensors, communication networks and software to improve urban services. Distributed
devices measure traffic flow, air quality, noise, temperature, water level or electricity
consumption. Data are transmitted to a local controller or platform, where algorithms detect trends,
identify abnormal conditions and select an appropriate action.

Street lighting provides a simple example. An adaptive system can combine a light sensor, motion
detector and timetable. Lamps dim when streets are empty and return to a safe level when activity is
detected. Energy savings are possible, but minimum illumination, response time and failure behaviour
must be specified and verified.

Traffic management is more complex. Cameras, induction loops and connected vehicles estimate queue
lengths. A controller can modify signal timing to reduce waiting time and emissions. Local
optimisation of one junction, however, may create congestion elsewhere. Engineers must therefore
model the complete network and define priorities for buses, emergency vehicles, cyclists and
pedestrians.

Smart buildings and grids can cooperate. Occupancy sensors adjust heating and ventilation, while a
district energy manager coordinates photovoltaic generation, batteries and electric-vehicle
charging. This requires devices from different manufacturers to exchange data. Interoperability and
open standards are essential. Edge computing can process urgent data near the sensors, reducing
latency and dependence on a distant platform.

Connectivity also creates risks. Location or camera data may reveal personal behaviour. A network
failure must not stop traffic lights, water pumps or emergency systems. Good design separates
critical functions, minimises stored data, applies cybersecurity measures and provides redundant
local fallback modes. A city is not truly smart if services remain inaccessible to people without
smartphones or digital skills.
\end{ecarticle}

\begin{tcolorbox}[title=\sffamily\bfseries Engineering Insight,colback=yellow!10,colframe=ECOrange]
The value of an urban dataset depends on the decision it supports. Collecting more data increases
communication, storage, cybersecurity and privacy costs; engineers should measure only what is needed
to operate, maintain and verify the service.
\end{tcolorbox}

\begin{engineeringfocus}
\textbf{Architecture of an adaptive lighting service}
\begin{center}
\resizebox{\linewidth}{!}{%
\begin{tikzpicture}[font=\scriptsize,>=Latex,node distance=4mm]
\node[draw=ECGreen,fill=ECGreen!10,rounded corners,align=center] (s) {motion and\\light sensors};
\node[draw=ECBlue,fill=ECLightBlue,rounded corners,align=center,right=of s] (n) {wireless\\network};
\node[draw=ECPurple,fill=ECPurple!10,rounded corners,align=center,right=of n] (c) {local\\controller};
\node[draw=ECOrange,fill=ECOrange!10,rounded corners,align=center,below=of c] (l) {LED street\\lamps};
\node[draw=ECRed,fill=ECRed!8,rounded corners,align=center,below=of n] (p) {supervision and\\maintenance platform};
\draw[->,thick] (s)--(n); \draw[->,thick] (n)--(c); \draw[->,thick] (c)--(l);
\draw[<->,thick,dashed] (n)--(p); \draw[->,thick,dashed] (l.west)-| (s.south);
\end{tikzpicture}%
}
\end{center}
The essential control loop is local: sensing, decision and lighting remain available even if the
central platform is disconnected.
\end{engineeringfocus}

\begin{keyfigures}
\begin{tabularx}{\linewidth}{@{}Xr@{}}
Connected devices & thousands per district\\
Useful lamp response time & a few seconds\\
Important design property & interoperability\\
Critical service requirement & local fallback\\
Main ethical concern & privacy
\end{tabularx}
\end{keyfigures}

\begin{tcolorbox}[title=\sffamily\bfseries Lighting power profile,colback=white,colframe=ECGreen]
\begin{center}
\begin{tikzpicture}
\begin{axis}[width=.96\linewidth,height=45mm,xlabel={Time of day (h)},ylabel={Rated power (\%)},xmin=0,xmax=24,ymin=0,ymax=110,grid=major,legend style={font=\scriptsize,at={(.5,-.28)},anchor=north,legend columns=2},ticklabel style={font=\scriptsize}]
\addplot[very thick,mark=*] table[x=hour,y=conventional,col sep=comma]{data/EC13/lighting_profile.csv};
\addlegendentry{Conventional}
\addplot[very thick,mark=square*] table[x=hour,y=adaptive,col sep=comma]{data/EC13/lighting_profile.csv};
\addlegendentry{Adaptive}
\end{axis}
\end{tikzpicture}
\end{center}
\footnotesize Illustrative daily profiles. The adaptive system lowers demand during quiet periods
while restoring a higher level when activity is detected.
\end{tcolorbox}

\begin{vocabulary}
\begin{tabularx}{\linewidth}{@{}lX@{}}
sensor network & réseau de capteurs\\
traffic flow & flux de circulation\\
street lighting & éclairage public\\
occupancy & occupation / présence\\
interoperability & interopérabilité\\
digital twin & jumeau numérique\\
fallback mode & mode de repli\\
outage & panne / interruption\\
privacy & vie privée\\
urban planning & urbanisme
\end{tabularx}
\end{vocabulary}

\begin{questions}
\item Describe the information and energy paths in the lighting system.
\item How can adaptive lighting save energy without reducing safety?
\item Why can local traffic optimisation be insufficient?
\item What does interoperability mean in this context?
\item Use the graph to identify the periods with the greatest potential savings.
\item Give three ways to make a connected urban service resilient.
\item Why should data minimisation be treated as an engineering requirement?
\end{questions}

\begin{applicationexercise}
A district operates 1,200 LED street lamps rated at 80 W. The conventional system runs at full power
for 12 hours each night. An adaptive strategy uses full power for 5 hours and 35\% power for the
remaining 7 hours.
\begin{enumerate}
\item Calculate the nightly energy consumption of both systems.
\item Determine the energy saved per night and the percentage reduction.
\item Estimate the annual saving over 365 nights.
\item State two effects not represented by this simplified energy calculation.
\end{enumerate}
\end{applicationexercise}

\begin{engineeringchallenge}
Design a smart energy-management system for a school district. Define requirements and performance
indicators; select sensors, communication links and controlled equipment; draw the information flow;
propose energy-saving rules; and describe local fallback modes. Include cybersecurity, maintenance,
data minimisation and access for users without smartphones. Present a justified architecture rather
than a list of connected devices.
\end{engineeringchallenge}

\begin{speaking}
Should a city use cameras and connected sensors to reduce congestion? Present a four-minute design
review covering one measurable benefit, one privacy risk, one resilience requirement and one
condition for public acceptance.
\end{speaking}
\end{multicols}

\subsection*{Sources}
\ECsource{European Commission -- Smart Cities}{https://commission.europa.eu/eu-regional-and-urban-development/topics/cities-and-urban-development/city-initiatives/smart-cities_en}
\ECsource{NIST -- Smart Connected Systems}{https://www.nist.gov/ctl/smart-connected-systems-division}
